\section{Анализ литературных источников и программных решений по теме дипломного проекта}

\subsection{Описание и анализ предметной области}

Сквозная аналитика для \eng{e-commerce} представляет собой подход к анализу эффективности интернет-продаж, при котором данные о переходах, пользовательских событиях, заказах, возвратах и выручке рассматриваются в едином контексте. Такой подход позволяет оценивать не только объем привлеченного трафика, но и вклад каждого источника в достижение бизнес-результата.

В отличие от базовой статистики посещаемости, сквозная аналитика рассматривает путь пользователя как последовательность связанных событий. В рамках интернет-магазина такой путь может начинаться с перехода из поисковой системы, рекламного объявления, социальной сети или рассылки. Далее пользователь просматривает страницу категории, открывает карточку товара, добавляет товар в корзину, оформляет заказ и выполняет оплату. Каждый этап может быть зафиксирован в виде отдельного события, а итоговая покупка связывается с источником привлечения. Благодаря этому становится возможным анализировать не только количество посещений, но и качество трафика.

Для \eng{e-commerce} проекта особую роль играет корректное сопоставление данных о переходах, событиях и заказах. Если рассматривать эти данные изолированно, можно получить искаженное представление об эффективности продвижения. Например, один источник может приводить большое количество посетителей, но иметь низкую конверсию в покупку. Другой источник может давать меньший объем трафика, но обеспечивать высокий средний чек и низкую долю возвратов. Поэтому система сквозной аналитики должна поддерживать расчет показателей в разрезе источников, периодов и этапов воронки продаж.

Для интернет-магазина важно понимать, какие каналы продвижения приводят к заявкам, покупкам и повторным обращениям. \eng{SEO}-специалист или аналитик должен сопоставлять данные о визитах, событиях на сайте, оформленных заказах, возвратах и рекламных расходах. Без специализированного программного средства эта работа часто выполняется вручную: данные собираются из разных таблиц и отчетов, проверяются на корректность и затем объединяются для расчета ключевых показателей.

Ручная обработка аналитических данных имеет несколько существенных недостатков. Во-первых, специалист тратит значительное время на сбор и приведение данных к единому виду. Во-вторых, при переносе информации из разных источников возрастает вероятность ошибок: дублирования строк, пропуска заказов, неверного указания периода или источника трафика. В-третьих, итоговая отчетность часто формируется в виде отдельных таблиц и документов, что затрудняет повторное использование результатов и сравнение показателей за разные периоды. В-четвертых, клиент не всегда имеет прозрачный доступ к текущему состоянию проекта и вынужден ожидать ручной подготовки отчета.

Сквозная аналитика позволяет снизить влияние перечисленных проблем за счет централизованного хранения данных и автоматизированного расчета показателей. При этом важна не только серверная обработка информации, но и удобный пользовательский интерфейс. Клиент должен понимать, какие данные необходимо передать системе, \eng{SEO}-специалист должен видеть статус обработки заявки и качество входных данных, а администратор должен иметь инструменты управления пользователями, проектами и демонстрационными наборами данных.

Разрабатываемое программное средство предназначено для автоматизации деятельности \eng{SEO}-специалиста при подключении и сопровождении сквозной аналитики \eng{e-commerce} проектов. Система обеспечивает регистрацию заявок клиентов, ведение проектов, ввод бизнес-данных в табличном виде, генерацию демонстрационного массива данных, проведение диагностики сайта, настройку источников трафика, расчет аналитических показателей, формирование рекомендаций и подготовку отчетов.

В рамках дипломного проекта система рассматривается как распределенное программное средство, состоящее из клиентской и серверной частей. Клиентская часть отвечает за отображение экранов, форм, таблиц, графиков и отчетов. Серверная часть обеспечивает хранение данных, проверку прав доступа, обработку заявок, расчет аналитических показателей и подготовку данных для отображения. Такое разделение соответствует распространенной практике разработки \eng{web}-систем и позволяет независимо описывать пользовательский интерфейс, прикладную логику и модель данных.

Основными данными, с которыми работает система, являются сведения о пользователях, заявках, проектах, источниках трафика, событиях, заказах, возвратах, отчетах и рекомендациях. Пользователь с ролью клиента создает заявку и передает исходную информацию о бизнесе. \eng{SEO}-специалист принимает заявку, выполняет диагностику, настраивает модель аналитики и анализирует рассчитанные показатели. Администратор управляет учетными записями, справочниками и демонстрационными данными, что позволяет использовать систему в учебной среде без подключения внешнего интернет-магазина.

В качестве результата работы системы используется \eng{dashboard} сквозной аналитики \eng{e-commerce} проекта, на котором отображаются ключевые показатели: переходы, заказы, выручка, конверсия, возвраты, средний чек и эффективность источников трафика. Пример такого \eng{dashboard} представлен на рисунке~\ref{fig:dashboard-ecommerce}.

\begin{figure}[H]
    \centering
    \diplomaimage{images/fig-1-1-dashboard-ecommerce.png}{68mm}
    \caption{Пример \eng{dashboard} сквозной аналитики \eng{e-commerce} проекта}
    \label{fig:dashboard-ecommerce}
\end{figure}

Сервис сквозной аналитики должен предоставлять пользователям единый интерфейс для контроля заявок, проектов, отчетов, источников трафика и аналитических задач. Клиент получает возможность передавать данные о бизнесе и просматривать результаты анализа, \eng{SEO}-специалист выполняет диагностику и формирует рекомендации, а администратор управляет пользователями, справочниками, тестовыми данными и резервным копированием.

Наличие единого интерфейса особенно важно для задач сопровождения нескольких проектов. В практической деятельности аналитический отдел или \eng{SEO}-агентство может одновременно работать с разными интернет-магазинами, отличающимися нишей, регионом, средним чеком, набором источников трафика и целями анализа. Если информация по каждому проекту хранится в отдельных файлах, возрастает сложность контроля сроков, статусов и результатов. Разрабатываемая система должна устранить эту проблему за счет единого пространства, в котором каждая заявка связана с клиентом, специалистом, проектом и последующими отчетами.

Для клиента ценность такого интерфейса заключается в прозрачности процесса. После подачи заявки клиент может видеть, на каком этапе находится обработка, какие данные уже переданы, какие показатели рассчитаны и какие рекомендации подготовлены. Для \eng{SEO}-специалиста интерфейс выступает рабочей панелью, где отображаются назначенные заявки, проекты, задачи, результаты проверки данных и аналитические выводы. Для администратора интерфейс обеспечивает контроль системных параметров, пользователей и тестовых данных.

Интерфейс обзора аналитической системы должен показывать общее состояние работы сервиса: количество активных проектов, клиентов, новых заявок, отчетов за период, а также динамику ключевых показателей. Пример интерфейса обзора аналитической системы представлен на рисунке~\ref{fig:analytics-overview}.

\begin{figure}[H]
    \centering
    \diplomaimage{images/fig-1-2-analytics-overview.png}{68mm}
    \caption{Интерфейс обзора аналитической системы}
    \label{fig:analytics-overview}
\end{figure}

Одним из начальных процессов системы является обработка заявки клиента на подключение сквозной аналитики. В заявке фиксируются сайт, ниша, цель аналитики, описание проблемы, выбранный \eng{SEO}-специалист, статус и дата создания. Такая структура позволяет централизовать входящие обращения и контролировать дальнейшую работу по каждому \eng{e-commerce} проекту.

Заявка выполняет роль исходной точки жизненного цикла проекта. На этом этапе клиент формулирует проблему, ради которой подключается аналитика: низкая конверсия, отсутствие понимания эффективности каналов, высокая доля возвратов, сложность подготовки отчетов или необходимость сравнить источники трафика. Дополнительно фиксируются сведения о сайте, нише, регионе, типе товаров, среднем чеке и ожидаемых результатах анализа. Наличие таких данных позволяет \eng{SEO}-специалисту быстрее оценить контекст проекта и определить первичные действия.

В системе предусмотрен выбор \eng{SEO}-специалиста из списка активных пользователей. Это позволяет уже на этапе подачи заявки связать клиента с ответственным исполнителем. При этом администратор может контролировать распределение нагрузки и при необходимости изменить назначение. Статусная модель заявки помогает отразить ее текущее состояние: новая заявка, обработка, запрос уточнения, принятие или отклонение. В дальнейшем принятая заявка может стать основанием для создания проекта, в котором будут храниться аналитические данные и отчеты.

Страница заявок клиентов используется для просмотра, фильтрации и изменения статусов обращений. Она помогает \eng{SEO}-специалисту видеть назначенные ему заявки, а администратору контролировать общую загрузку специалистов. Пример страницы заявок клиентов представлен на рисунке~\ref{fig:client-applications}.

\begin{figure}[H]
    \centering
    \diplomaimage{images/fig-1-3-client-applications.png}{68mm}
    \caption{Страница заявок клиентов}
    \label{fig:client-applications}
\end{figure}

Таким образом, предметная область дипломного проекта связана с автоматизацией процессов сбора, проверки, анализа и представления данных \eng{e-commerce} проектов. Основная задача программного средства заключается в сокращении ручного труда \eng{SEO}-специалиста, повышении прозрачности эффективности источников трафика и формировании понятной отчетности для клиента.

Ключевыми аналитическими показателями предметной области являются количество переходов, количество событий, число заказов, выручка, конверсия, средний чек, сумма возвратов, доля возвратов, стоимость привлечения заказа и эффективность источников трафика. Для учебного программного средства достаточно реализовать расчет этих показателей на основе табличных данных, вводимых пользователем или сформированных администратором. Такой подход не требует внешней интеграции, но позволяет показать принцип работы сквозной аналитики на реалистичной модели данных.

При проектировании такой системы необходимо учитывать качество входных данных. Если в таблицах отсутствует дата события, не указан источник трафика, сумма заказа записана в некорректном формате или возврат не связан с заказом, итоговые аналитические показатели будут неточными. Поэтому в предметной области важна не только визуализация результата, но и предварительная проверка данных. Система должна выявлять типовые ошибки, сообщать пользователю о проблемах и не допускать расчета показателей на заведомо некорректном наборе записей.

Дополнительной особенностью является необходимость хранения истории аналитических работ. Для клиента важно видеть не только текущий отчет, но и динамику изменений: как изменялась конверсия, какие рекомендации были сформированы ранее, какие источники трафика улучшили показатели, а какие остались проблемными. Для \eng{SEO}-специалиста история проекта помогает обосновывать выводы и отслеживать влияние принятых решений. Поэтому разрабатываемое программное средство должно рассматриваться как система сопровождения проекта, а не как одноразовый калькулятор показателей.

Следовательно, разрабатываемая система должна объединять признаки учетной системы, аналитического сервиса и средства визуализации. Учетная часть отвечает за заявки, пользователей, проекты и статусы. Аналитическая часть рассчитывает показатели и формирует выводы. Визуальная часть представляет результат в понятной форме: карточки показателей, графики, таблицы, отчеты и рекомендации. Совместная работа этих компонентов обеспечивает достижение цели дипломного проекта и формирует основу для дальнейшего проектирования требований.

\subsection{Обзор функциональности аналогов программного средства}

На рынке представлено несколько классов решений, которые частично закрывают задачи \eng{web}-аналитики, визуализации данных и сквозной аналитики. К ним относятся системы \eng{web}-аналитики, \eng{BI}-платформы, сервисы построения \eng{dashboard} и специализированные решения для анализа маркетинговой эффективности. Для оценки текущей ситуации рассмотрим наиболее известные программные продукты, применяемые при анализе \eng{e-commerce} проектов. Конкуренты на рынке программных решений для сквозной аналитики \eng{e-commerce} представлены в таблице~\ref{tab:analytics-competitors}.

При выборе аналогов учитывались решения, которые используются для сбора данных о поведении пользователей, построения отчетов, визуализации показателей и оценки маркетинговой эффективности. Рассматриваемые программные продукты отличаются назначением. Одни из них ориентированы на регистрацию событий и анализ посещаемости, другие применяются для построения интерактивной отчетности, третьи специализируются на сквозной аналитике и оценке окупаемости рекламных каналов. Для дипломного проекта важно сопоставить эти решения с задачами разрабатываемой системы и определить, какие функции должны быть отражены в учебном программном средстве.

Системы \eng{web}-аналитики, такие как \eng{Google Analytics} и Яндекс Метрика, позволяют отслеживать посещения сайта, источники переходов, цели, события и аудитории. Они хорошо подходят для анализа поведения пользователей, но не всегда закрывают задачу управления заявками клиентов, назначения специалистов и подготовки внутренней отчетности по проектам. Кроме того, такие системы предполагают настройку счетчиков и интеграций с реальным сайтом, что выходит за рамки образовательного проекта, в котором внешний интернет-магазин не разрабатывается как отдельный компонент.

\eng{BI}-платформы, например \eng{Looker Studio} и \eng{Power BI}, предоставляют развитые средства построения отчетов и визуализации данных. Их сильной стороной является гибкость представления информации, подключение разных источников и создание интерактивных панелей. Однако такие инструменты в большей степени являются универсальными конструкторами отчетности, чем специализированными системами сопровождения сквозной аналитики. В них обычно отсутствует предметная логика заявок, ролей клиента и \eng{SEO}-специалиста, проверки качества табличных данных и генерации демонстрационного массива для учебного сценария.

Специализированные сервисы сквозной аналитики, такие как \eng{Roistat} и \eng{Calltouch}, ближе к теме дипломного проекта, поскольку ориентированы на оценку эффективности каналов привлечения и связку маркетинговых данных с заявками или продажами. При этом такие сервисы являются готовыми коммерческими продуктами с большим количеством интеграций и настроек, часть которых не требуется в учебной реализации. Разрабатываемое программное средство должно быть проще по масштабу, но должно явно демонстрировать основные процессы: заявку, проект, ввод данных, расчет показателей, \eng{dashboard}, рекомендации и отчет.

\begin{table}[H]
    \caption{Конкуренты на рынке программных решений для сквозной аналитики \eng{e-commerce}}
    \label{tab:analytics-competitors}
    {\small
    \setlength{\tabcolsep}{4pt}
    \begin{tabularx}{\textwidth}{|P{0.24\textwidth}|L|P{0.24\textwidth}|}
        \hline
        Конкурент & Возможности & Ссылка на решение \\
        \hline
        \eng{Google Analytics} & Анализ \eng{web}-трафика, событий, конверсий, источников переходов и аудиторий & \url{analytics.google.com} \\
        \hline
        Яндекс Метрика & \eng{Web}-аналитика, цели, сегменты, отчеты по источникам трафика, вебвизор & \url{metrika.yandex.ru} \\
        \hline
        \eng{Roistat} & Сквозная аналитика, отчеты по рекламным каналам, расчет эффективности маркетинга & \url{roistat.com} \\
        \hline
        \eng{Looker Studio} & Визуализация данных, построение \eng{dashboard}, подключение внешних источников & \url{lookerstudio.google.com} \\
        \hline
        \eng{Power BI} & \eng{BI}-аналитика, интерактивные отчеты, обработка и визуализация данных & \url{powerbi.microsoft.com} \\
        \hline
        \eng{Calltouch} & Сквозная аналитика, коллтрекинг, отчеты по каналам привлечения и лидам & \url{calltouch.ru} \\
        \hline
    \end{tabularx}
    }
\end{table}

Рассмотренные решения обладают развитой функциональностью, однако каждое из них ориентировано на собственный набор задач. Системы \eng{web}-аналитики хорошо подходят для анализа посещаемости и поведения пользователей, \eng{BI}-платформы удобны для построения отчетности, а специализированные сервисы сквозной аналитики требуют настройки источников данных и интеграций. Для учебного проекта важно разработать программное средство, которое демонстрирует полный цикл обработки данных \eng{e-commerce} проекта без подключения внешнего интернет-магазина или стороннего \eng{API}.

Сравнение аналогов показывает, что для целей дипломного проекта недостаточно просто воспроизвести аналитический экран. Требуется программное средство, в котором аналитика связана с организационным процессом. Клиент должен иметь возможность подать заявку, выбрать специалиста и передать данные. \eng{SEO}-специалист должен принять заявку, проверить сведения, настроить модель аналитики и подготовить рекомендации. Администратор должен управлять пользователями и демонстрационными данными. Именно такая связка процессов отличает разрабатываемую систему от универсальных инструментов отчетности.

Также следует учитывать, что большинство существующих решений ориентированы на работу с внешними источниками данных. Для реальной промышленной системы это является преимуществом, но в учебном проекте такая зависимость усложняет демонстрацию. Если работоспособность дипломного проекта зависит от стороннего сервиса, учетной записи или внешнего интернет-магазина, проверка системы становится менее управляемой. Поэтому в разрабатываемом программном средстве используются ручной ввод табличных данных и генерация демонстрационного массива. Это делает систему автономной и позволяет проверять ее без внешних интеграций.

Сравнение функциональных возможностей решений сквозной аналитики приведено в таблице~\ref{tab:analytics-features}.

\begin{table}[H]
    \caption{Функциональные возможности решений сквозной аналитики}
    \label{tab:analytics-features}
    {\small
    \setlength{\tabcolsep}{3pt}
    \begin{tabularx}{\textwidth}{|P{0.23\textwidth}|Y|Y|Y|Y|Y|Y|}
        \hline
        Функциональная возможность & \eng{Google Analytics} & Яндекс Метрика & \eng{Roistat} & \eng{Looker Studio} & \eng{Power BI} & \eng{Calltouch} \\
        \hline
        \eng{Dashboard} с \eng{KPI} & Да & Да & Да & Да & Да & Да \\
        \hline
        Анализ источников трафика & Да & Да & Да & Частично & Частично & Да \\
        \hline
        Анализ воронки продаж & Да & Да & Да & Частично & Частично & Частично \\
        \hline
        Импорт табличных данных & Частично & Частично & Частично & Да & Да & Частично \\
        \hline
        Управление заявками клиентов & Нет & Нет & Частично & Нет & Нет & Частично \\
        \hline
        Ролевая модель для клиента, \eng{SEO}-специалиста и администратора & Частично & Частично & Частично & Частично & Частично & Частично \\
        \hline
        Формирование рекомендаций & Нет & Нет & Частично & Нет & Нет & Частично \\
        \hline
        Генерация демонстрационных данных & Нет & Нет & Нет & Нет & Нет & Нет \\
        \hline
    \end{tabularx}
    }
\end{table}

Анализ аналогов показывает, что существующие решения позволяют решать отдельные задачи \eng{e-commerce} аналитики, но не всегда предоставляют простой учебный сценарий полного цикла: от заявки клиента и ввода данных до расчета \eng{KPI}, формирования \eng{dashboard}, рекомендаций и отчета. Разрабатываемое программное средство должно объединить эти этапы в единую \eng{web}-систему и обеспечить демонстрацию сквозной аналитики на тестовых или вручную введенных данных.

По результатам сравнения можно выделить несколько требований к разрабатываемому программному средству. Во-первых, система должна быть предметно ориентированной, то есть использовать сущности, характерные для сопровождения аналитики интернет-магазинов: заявки, проекты, источники трафика, события, заказы, возвраты, отчеты и рекомендации. Во-вторых, интерфейс должен быть понятным для разных ролей пользователей и не ограничиваться только отображением графиков. В-третьих, аналитическая часть должна рассчитывать основные показатели по данным проекта и представлять их в виде, пригодном для принятия решений. В-четвертых, система должна быть автономной и воспроизводимой в учебной среде.

С точки зрения функциональности разрабатываемая система занимает промежуточное положение между сервисом заявок, аналитической платформой и системой отчетности. От сервиса заявок она наследует регистрацию обращений клиентов, статусы и назначение исполнителя. От аналитической платформы -- расчет показателей, анализ источников трафика и построение \eng{dashboard}. От системы отчетности -- формирование итоговых выводов и рекомендаций для клиента. Такое сочетание делает проект целостным и соответствует цели дипломной работы.

Важным отличием от рассмотренных аналогов является ориентация на демонстрационный сценарий. Система не должна выдавать себя за коммерчески внедренный продукт и не должна описываться как фактически продаваемый сервис. Она предназначена для показа архитектуры, пользовательских сценариев, модели данных и логики сквозной аналитики. Поэтому экономические преимущества в дальнейших разделах должны рассматриваться как условная оценка возможного эффекта от автоматизации, а не как данные о реальных продажах или коммерческой эксплуатации.

Кроме функциональных отличий, для дипломного проекта важны критерии простоты развертывания и воспроизводимости. Готовые аналитические платформы часто требуют настройки аккаунтов, подключения счетчиков, импорта из рекламных кабинетов или подготовки сложных источников данных. В учебной системе целесообразно сделать сценарий более контролируемым: пользователь вводит данные через интерфейс, администратор при необходимости генерирует тестовый набор, а серверная часть рассчитывает показатели на основе локальной базы данных. Такой подход позволяет продемонстрировать работу системы при защите дипломного проекта без зависимости от внешних учетных записей.

Также необходимо учитывать различие между универсальными средствами визуализации и предметно-ориентированным программным средством. Универсальная платформа может построить график по загруженной таблице, однако она не знает, что такое заявка на подключение аналитики, статус проекта, роль \eng{SEO}-специалиста, возврат заказа или рекомендация по источнику трафика. В разрабатываемой системе эти понятия должны быть встроены в модель данных и пользовательские сценарии. Благодаря этому интерфейс становится ближе к реальной деятельности аналитического отдела, а не просто к отображению произвольных диаграмм.

Проведенный обзор позволяет сформулировать общее направление разработки. Программное средство должно сохранять простоту учебного проекта, но при этом отражать основные элементы профессиональной работы со сквозной аналитикой: сбор исходных данных, контроль качества, расчет показателей, интерпретацию результатов и подготовку отчета. Именно поэтому в дальнейшем при моделировании предметной области необходимо описать не только функциональные экраны, но и бизнес-процессы, роли пользователей, ограничения доступа и требования к данным.

С учетом рассмотренных аналогов можно сделать вывод, что разрабатываемая система должна быть ориентирована не на конкуренцию с промышленными аналитическими платформами, а на демонстрацию полного и понятного учебного цикла. В промышленной системе значительная часть сложности связана с интеграциями, обработкой больших объемов данных и поддержкой множества внешних сервисов. В дипломном проекте приоритет имеет другое: показать, какие сущности нужны для предметной области, как пользователь проходит основной сценарий, какие данные сохраняются в базе, как рассчитываются показатели и как результат становится доступен клиенту.

Такой подход позволяет сохранить баланс между реалистичностью и выполнимостью проекта. Если система будет слишком упрощенной, она не покажет специфику сквозной аналитики и сведется к обычному набору таблиц. Если же система будет перегружена внешними интеграциями, появится риск выйти за рамки дипломного проекта и потерять контроль над демонстрацией. Поэтому выбранный вариант основан на автономной обработке данных: клиент или администратор формирует набор переходов, событий, заказов и возвратов, после чего серверная часть рассчитывает показатели и подготавливает данные для интерфейса.

Для последующего проектирования важно, что анализ аналогов выявил не только полезные функции, но и ограничения существующих решений применительно к учебной задаче. В разрабатываемой системе должны быть сохранены наиболее значимые элементы: аналитическая панель, оценка источников трафика, отчеты, роли пользователей и работа с проектами. При этом должны быть исключены лишние для учебной реализации зависимости от внешних кабинетов, платных интеграций и сторонних интернет-магазинов. Это делает программное средство самодостаточным и позволяет в последующих разделах последовательно описать требования, архитектуру, модель данных и пользовательские сценарии.

Для пользователя итоговая ценность системы выражается в том, что результаты анализа становятся доступными в едином интерфейсе, а не в виде разрозненных файлов. Клиент получает понятную картину по проекту: сколько было переходов, сколько оформлено заказов, какая выручка получена за выбранный период, какова конверсия и какие источники требуют внимания. \eng{SEO}-специалист получает инструмент для системной работы с данными: он может проверить полноту сведений, сопоставить источники, выделить слабые этапы воронки и подготовить рекомендации. Администратор получает возможность поддерживать демонстрационную среду, управлять пользователями и обеспечивать воспроизводимость сценариев.

С методической точки зрения такая система удобна для дипломного проекта, потому что каждый ее элемент можно обосновать и проверить. Заявки подтверждают наличие организационного процесса, проекты связывают клиента и специалиста, источники трафика задают аналитическое измерение, события и заказы формируют исходные данные, отчеты и рекомендации фиксируют результат работы. Благодаря этому пояснительная записка может последовательно раскрывать переход от анализа предметной области к требованиям, от требований к архитектуре, от архитектуры к модели данных и пользовательскому интерфейсу.

Таким образом, обзор аналогов выполняет не только сравнительную функцию, но и задает границы будущей разработки. В проект не включается разработка внешнего интернет-магазина, подключение реальных рекламных кабинетов и коммерческая модель фактических продаж. Вместо этого внимание сосредоточено на программном средстве, которое показывает основной смысл сквозной аналитики: объединение данных о трафике, действиях пользователей, заказах и возвратах в единую систему показателей.

Таким образом, анализ программных решений подтверждает необходимость разработки учебного программного средства, которое объединяет управление заявками, сопровождение проектов, ввод табличных данных, расчет аналитических показателей и подготовку отчетности. Рассмотренные аналоги показывают существующие подходы к аналитике и визуализации, однако не закрывают полностью учебный сценарий, выбранный для дипломного проекта. Это обосновывает дальнейший переход к моделированию предметной области и формированию требований к системе.
